\section*{Bloque C. Casos con implicancias regulatorias y de política}

\noindent\textit{Preguntas 66 a 90. Aplican el equilibrio general, la eficiencia de Pareto, la caja de Edgeworth y los teoremas del bienestar a la asignación de agua, los derechos de agua, las externalidades, el monopolio natural y la evaluación de políticas tarifarias.}

\noindent\textit{Para las preguntas 66 a 69, considere dos usuarios de agua, A y B, con un total de 100 unidades. Los beneficios marginales son $MBA=100-q_A$ y $MBB=80-q_B$, y los beneficios totales son las áreas bajo estas curvas.}

\q{La asignación eficiente de agua para el usuario A es:}
\begin{opts}
  \item 50
  \item 60
  \item 70
  \item 80
\end{opts}

\q{La asignación eficiente de agua para el usuario B es:}
\begin{opts}
  \item 30
  \item 40
  \item 50
  \item 60
\end{opts}

\q{Si la asignación inicial es A$=70$ y B$=30$, la ganancia total de beneficio al pasar a la asignación eficiente es:}
\begin{opts}
  \item 50
  \item 75
  \item 100
  \item 200
\end{opts}

\q{En una caja de Edgeworth con 100 unidades de agua y 200 unidades del otro bien, si A tiene $(60,80)$, la cantidad del otro bien que tiene B es:}
\begin{opts}
  \item 80
  \item 120
  \item 100
  \item 60
\end{opts}

\q{Si en una asignación de agua $MRS_A=3$ y $MRS_B=1$, entonces:}
\begin{opts}
  \item ambos usuarios pueden ganar con el intercambio.
  \item no existen ganancias de intercambio.
  \item solo pueden ganar si los precios son iguales.
  \item solo el regulador puede mejorar la asignación.
\end{opts}

\q{En un óptimo de Pareto interior de asignación de agua entre dos usuarios se cumple:}
\begin{opts}
  \item $MRS_A>MRS_B$.
  \item $MRS_A=MRS_B$.
  \item $MRS=0$.
  \item $MRS=MRT$ únicamente.
\end{opts}

\q{Si una asignación de agua se encuentra sobre la curva de contrato, entonces:}
\begin{opts}
  \item no es eficiente.
  \item es necesariamente equitativa.
  \item maximiza la producción física.
  \item es Pareto eficiente.
\end{opts}

\q{Si la asignación inicial es A$=80$ y B$=20$, y la asignación eficiente es A$=60$ y B$=40$, la transferencia de agua de A hacia B es:}
\begin{opts}
  \item 20
  \item 10
  \item 30
  \item 40
\end{opts}

\q{La reasignación de 20 unidades desde A hacia B aumenta el beneficio total. Esto implica que:}
\begin{opts}
  \item existen ganancias potenciales de bienestar; con compensación puede ser una mejora de Pareto.
  \item A mejora con seguridad aunque no reciba compensación.
  \item B necesariamente pierde con la reasignación.
  \item no existen ganancias de bienestar.
\end{opts}

\q{Si los derechos de agua están bien definidos y los costos de transacción son cero, el teorema de Coase implica que:}
\begin{opts}
  \item solo el gobierno puede asignar el agua eficientemente.
  \item la asignación eficiente requiere repartir el agua por igual.
  \item no habrá intercambio de derechos.
  \item se puede alcanzar una asignación eficiente con independencia de la asignación inicial de derechos.
\end{opts}

\q{Para que el Primer Teorema del Bienestar se aplique a un mercado de agua se requiere, entre otros supuestos:}
\begin{opts}
  \item existencia de monopolio y barreras a la entrada.
  \item ausencia de externalidades, mercados completos y agentes precio-aceptantes.
  \item externalidades negativas y bienes públicos.
  \item precios máximos fijados por el regulador.
\end{opts}

\q{Si la extracción de agua genera una externalidad negativa, el equilibrio competitivo:}
\begin{opts}
  \item es siempre Pareto eficiente.
  \item es necesariamente equitativo.
  \item no es Pareto eficiente.
  \item tiene exceso de demanda cero.
\end{opts}

\q{Un impuesto pigouviano sobre la extracción de agua debe fijarse igual a:}
\begin{opts}
  \item cero.
  \item el costo medio total.
  \item el costo total del servicio.
  \item el costo externo marginal.
\end{opts}

\q{El Segundo Teorema del Bienestar aplicado al agua implica que:}
\begin{opts}
  \item todo equilibrio competitivo es equitativo.
  \item una asignación eficiente puede sostenerse con precios y transferencias lump-sum, bajo ciertos supuestos.
  \item la eficiencia exige un monopolio estatal.
  \item se debe fijar un precio máximo.
\end{opts}

\q{Para proteger a los hogares vulnerables sin distorsionar las señales de precio, la opción más adecuada es:}
\begin{opts}
  \item un precio máximo general por debajo del equilibrio.
  \item un precio mínimo.
  \item una cuota de consumo.
  \item una transferencia lump-sum focalizada.
\end{opts}

\q{En un monopolio natural de agua, fijar el precio igual al costo marginal genera:}
\begin{opts}
  \item beneficios extraordinarios.
  \item beneficio económico nulo.
  \item ausencia de pérdidas.
  \item pérdidas que requieren un subsidio o una tarifa en dos partes.
\end{opts}

\q{La protección de una cuenca como bien público se caracteriza por:}
\begin{opts}
  \item ser rival y excluible.
  \item ser provista en exceso por el mercado.
  \item no ser rival ni excluible, por lo que el mercado tiende a subproveerla.
  \item ser provista eficientemente por una empresa competitiva.
\end{opts}

\q{La condición de Samuelson para la provisión eficiente de un bien público es:}
\begin{opts}
  \item $\sum MRS=MRT$.
  \item $MRS=0$.
  \item $MRS=MRT$.
  \item $\sum MRS=0$.
\end{opts}

\q{El criterio de Kaldor-Hicks sostiene que una política es deseable si:}
\begin{opts}
  \item nadie empeora.
  \item los ganadores podrían compensar a los perdedores y aun así mejorar.
  \item todos reciben la misma utilidad.
  \item el gobierno fija los precios.
\end{opts}

\q{Si los ganadores compensan efectivamente a los perdedores y nadie queda peor, el resultado es:}
\begin{opts}
  \item una mejora de Pareto.
  \item solo una mejora potencial de Kaldor-Hicks.
  \item una asignación ineficiente.
  \item un precio máximo.
\end{opts}

\q{Una asignación Pareto eficiente:}
\begin{opts}
  \item es siempre equitativa.
  \item puede ser eficiente pero inequitativa, por lo que se requiere un juicio distributivo.
  \item siempre iguala las utilidades.
  \item elimina la necesidad de redistribución.
\end{opts}

\q{Una función de bienestar social rawlsiana se concentra en:}
\begin{opts}
  \item la suma de las utilidades.
  \item el criterio de Kaldor-Hicks.
  \item la utilidad del individuo peor situado.
  \item la eficiencia productiva.
\end{opts}

\q{Una función de bienestar social utilitarista busca maximizar:}
\begin{opts}
  \item la suma de las utilidades individuales.
  \item la utilidad del peor situado.
  \item la desigualdad.
  \item el precio del agua.
\end{opts}

\q{Al evaluar una reforma tarifaria de agua en equilibrio general, deben considerarse:}
\begin{opts}
  \item únicamente el mercado de agua.
  \item únicamente a los consumidores.
  \item únicamente a los productores.
  \item los efectos en otros mercados, como alimentos, energía y trabajo.
\end{opts}

\q{La evaluación de bienestar de una política de agua debe incluir:}
\begin{opts}
  \item solo el excedente del consumidor.
  \item las externalidades y los efectos de equilibrio general, además de los excedentes.
  \item solo el excedente del productor.
  \item solo la recaudación fiscal.
\end{opts}
